%%%%%%%%%%%%%%%%%%%%%%%%%%%%%%%%%%%%%%%%%%
% Short Sectioned Assignment LaTeX Template Version 1.0 (5/5/12)
% This template has been downloaded from: http://www.LaTeXTemplates.com
% Original author:  Frits Wenneker (http://www.howtotex.com)
% License: CC BY-NC-SA 3.0 (http://creativecommons.org/licenses/by-nc-sa/3.0/)
%%%%%%%%%%%%%%%%%%%%%%%%%%%%%%%%%%%%%%%%%

%----------------------------------------------------------------------------------------
%	PACKAGES AND OTHER DOCUMENT CONFIGURATIONS
%----------------------------------------------------------------------------------------

\documentclass[paper=a4, fontsize=11pt]{scrartcl} % A4 paper and 11pt font size

% ---- Entrada y salida de texto -----

\usepackage[T1]{fontenc} % Use 8-bit encoding that has 256 glyphs
\usepackage[utf8]{inputenc}
%\usepackage{fourier} % Use the Adobe Utopia font for the document - comment this line to return to the LaTeX default

% ---- Idioma --------

\usepackage[spanish, es-tabla]{babel} % Selecciona el español para palabras introducidas automáticamente, p.ej. "septiembre" en la fecha y especifica que se use la palabra Tabla en vez de Cuadro

% ---- Otros paquetes ----

\usepackage{url} % ,href} %para incluir URLs e hipervínculos dentro del texto (aunque hay que instalar href)
\usepackage{amsmath,amsfonts,amsthm} % Math packages
%\usepackage{graphics,graphicx, floatrow} %para incluir imágenes y notas en las imágenes
\usepackage{graphics,graphicx, float} %para incluir imágenes y colocarlas

% Para hacer tablas comlejas
%\usepackage{multirow}
%\usepackage{threeparttable}

%\usepackage{sectsty} % Allows customizing section commands
%\allsectionsfont{\centering \normalfont\scshape} % Make all sections centered, the default font and small caps

\usepackage{fancyhdr} % Custom headers and footers
\pagestyle{fancyplain} % Makes all pages in the document conform to the custom headers and footers
\fancyhead{} % No page header - if you want one, create it in the same way as the footers below
\fancyfoot[L]{} % Empty left footer
\fancyfoot[C]{} % Empty center footer
\fancyfoot[R]{\thepage} % Page numbering for right footer
\renewcommand{\headrulewidth}{0pt} % Remove header underlines
\renewcommand{\footrulewidth}{0pt} % Remove footer underlines
\setlength{\headheight}{13.6pt} % Customize the height of the header

\numberwithin{equation}{section} % Number equations within sections (i.e. 1.1, 1.2, 2.1, 2.2 instead of 1, 2, 3, 4)
\numberwithin{figure}{section} % Number figures within sections (i.e. 1.1, 1.2, 2.1, 2.2 instead of 1, 2, 3, 4)
\numberwithin{table}{section} % Number tables within sections (i.e. 1.1, 1.2, 2.1, 2.2 instead of 1, 2, 3, 4)

\setlength\parindent{0pt} % Removes all indentation from paragraphs - comment this line for an assignment with lots of text

\newcommand{\horrule}[1]{\rule{\linewidth}{#1}} % Create horizontal rule command with 1 argument of height






%----------------------------------------------------------------------------------------
%	TÍTULO Y DATOS DEL ALUMNO
%----------------------------------------------------------------------------------------

\title{	
\normalfont \normalsize 
\textsc{\textbf{Grundlagen der Wissensverarbeitung} \\ Computer Science \\ Universität Hamburg} \\ [25pt] % Your university, school and/or department name(s)
~\\
~\\
~\\
\horrule{0.5pt} \\[0.4cm] % Thin top horizontal rule
\Huge Tutorial 6: Search and Parsing \\ % The assignment title
\horrule{2pt} \\[0.5cm] % Thick bottom horizontal rule
~\\
~\\
}

\author{Rafael Ruz Gómez\\Miguel Robles Urquiza} % Nombre y apellidos

\date{\normalsize 27 November 2017} % Incluye la fecha actual

%----------------------------------------------------------------------------------------
% DOCUMENTO
%----------------------------------------------------------------------------------------

\begin{document}

\maketitle % Muestra el Título

\begin{figure}
	\centering
	\includegraphics[scale=0.8]{logo_uni_hamburg.png}
\end{figure}

\newpage %inserta un salto de página




%----------------------------------------------------------------------------------------
%	Question 1
%----------------------------------------------------------------------------------------

\section*{Exercise 1}

\subsection*{(a) By what operations is the input transformed into the output? What do these operations do?}

\begin{itemize}
	\item Left-Arc: Adds an arc from the next input token to the node on top of the stack and reduces from the stack(pops).
	\item Right-Arc: Adds an arc from the node on top of the stack to the next input token and pushes the next input token into the stack.
	\item Reduce: Pops from the stack.
	\item Shift: Push the next input item into the stack.
\end{itemize}

\subsection*{(b) When does the parsing algorithm terminate?}

It terminates when the list of remaining input tokens is empty, i.e. when it reaches a configuration <S,nil,A> where S represents the stack, the second element represents the empty input, and A represents the arc relation.

\subsection*{(c) Describe the formal properties of a dependency tree as defined in the paper.}

\begin{itemize}
	\item Single head: If there are 2 nodes having an arc to the same node, then those 2 nodes are the same node, i.e. every node has got 0 or 1 parent node, but no more than 1.
	\item Acyclic: For every 2 differents nodes n and n', there's never an arc from n to n' and at the same time one path(one or more arcs consecutively) from n' to n.
	\item Connected: For every 2 different nodes n and n', there's always a path between them (from n to n' or from n' to n).
	\item Projective: For every two nodes n and n' having an arc from one to the other, if there's another node n'' between them ( n < n'' < n' means POS(n) < POS(n'') < POS(n') ), then there has to be a path from n or from n' to n''
\end{itemize}

\subsection*{For each property: Give an example dependency tree that violates the property. Note: We ask for a tree, not a sentence! Do not try to find a matching sentence to your trees, it will only distract you.}

\begin{itemize}
	\item Single head: Three nodes n0,n1 and n2. Arcs: n0->n1 , n2->n1
	\item Acyclic: Two nodes n0 and n1. Arcs: n0->n1 , n1->n0
	\item Connected: Two nodes n0 and n1. Arcs: none
	\item Projective: Three nodes n0,n1 and n2. Note that the list of input tokens are n0=(0,w0), n1=(1,w1) and n2=(2,w2) where the pair is consisting of a position and a word form. Then the list of arcs would be: n0->n2
\end{itemize}



%----------------------------------------------------------------------------------------
%	Question 2
%----------------------------------------------------------------------------------------
\section*{Exercise 2}


\subsection*{Try to use the proposed parser actions to produce the tree depicted in Figure 1. Write down the steps and the intermediate states.}

\begin{enumerate}
	\item Initialization . <nil, Der Mann isst eine Giraffe, $\emptyset$>
	\item Shift . <Der, Mann isst eine Giraffe, $\emptyset$>
	\item Left-Arc . <nil, Mann isst eine Giraffe, {(Mann,Der)}>
	\item Shift . <Mann, isst eine Giraffe, {(Mann,Der)}>
	\item Left-Arc . <nil, isst eine Giraffe, {(Mann,Der),(isst,Mann)}>
	\item Shift . <isst, eine Giraffe, {(Mann,Der),(isst,Mann)}>
	\item Shift . <eine isst, Giraffe, {(Mann,Der),(isst,Mann)}>
	\item Left-Arc . <isst, Giraffe, {(Mann,Der),(isst,Mann), (Giraffe,eine)}>
	\item Right-Arc . <Giraffe isst, \textbf{nil}, {(Mann,Der),(isst,Mann), (Giraffe,eine), (isst,Giraffe)}>
	
\end{enumerate}





%----------------------------------------------------------------------------------------
%	Question 3
%----------------------------------------------------------------------------------------
\section*{Exercise 3}

\subsection*{The basic idea laid out in the paper by Nivre is still in use in state-of-the-art parsers such as the one coined “Parsey McParseface” by Google. While the actions are still the same, the mechanics of selecting actions are quite different. The actions are not ordered using a fixed precedence and their applicability is not restricted by lexical rules as they are introduced in the paper. Instead, search is used. If you now view parsing using the proposed actions as a search problem:}

\subsubsection*{• What are the search states?}
The search states are nodes composed by the set <S,I,A> as described before (S stack, I input tokens, A arcs). 

\subsubsection*{• What is the start state?}
The start state is a node composed by <nil, I, $\emptyset$>

\subsubsection*{• What are the end states?}
The end states are those who have no more input tokens remaining, i.e. <S, nil, A> where $A\neq\emptyset$

\subsubsection*{• What are the state transitions?}
The transitions are the same described in the paper :
\begin{itemize}
	\item Left-Arc
	\item Right-Arc
	\item Reduce
	\item Shift
\end{itemize}

\subsubsection*{• Can the search space be created before parsing starts?}
Yes, because the algorithm is deterministic. There are two kinds of transitions: Those that pops an element from the stack, and those that push an element into the stack(and then delete an element from the input). When the stack is empty, the only option is to make a “push” transition, so in the end, we will reach an end state, and so on we can create the search space before parsing starts.

\subsubsection*{• What is the advantage of the proposed algorithm in contrast to simply trying to find a good dependency tree by enumerating all possible trees and selecting the best one from them?}
The efficiency in time and even more in space.

\subsubsection*{• For the search strategies discussed so far: are they a good fit for this search problem and why (not)?}
It depends on the length of the input, but in fact depth-first looks like we are applying the algorithm with the priority of transitions described in the paper, so it's not a good idea to create the search space and then just do the same proposed in the paper.\\

Breadth-first has the same problem as always, the efficiency in time and space.\\

A* looks like the best option, but on the other hand I find not easy to define a good heuristic function.

\subsubsection*{• How would you design a parser using the parser actions together with an appropriate search procedure?}
I would define the search space as described above and then choose A* as the search algorithm.


\newpage
%------------------------------------------------

%\bibliography{citas} %archivo citas.bib que contiene las entradas 
%\bibliographystyle{plain} % hay varias formas de citar

\end{document}

